\documentclass[11pt]{article}

\usepackage[margin=0.72in]{geometry}
\usepackage[T1]{fontenc}
\usepackage[utf8]{inputenc}
\usepackage{lmodern}
\usepackage[table]{xcolor}
\usepackage{array}
\usepackage{longtable}
\usepackage{tabularx}
\usepackage{enumitem}
\usepackage{xurl}
\usepackage{hyperref}
\usepackage{fancyhdr}
\usepackage{titlesec}
\usepackage{microtype}
\usepackage{booktabs}
\usepackage{ragged2e}
\usepackage{lastpage}

\definecolor{ink}{HTML}{1F2937}
\definecolor{muted}{HTML}{4B5563}
\definecolor{accent}{HTML}{0F766E}
\definecolor{accentlight}{HTML}{E6F4F1}
\definecolor{rulegray}{HTML}{D1D5DB}
\definecolor{headergray}{HTML}{F3F4F6}

\hypersetup{
  colorlinks=true,
  urlcolor=accent,
  linkcolor=accent,
  citecolor=accent,
  pdftitle={Open-Source Financial Tools: Desktop and Self-Hosted Options},
  pdfauthor={Jordan Suber},
  pdfsubject={Citation-clean comparison of open-source personal finance, budgeting, accounting, and portfolio tools},
  pdfkeywords={open-source finance, GnuCash, HomeBank, Money Manager Ex, Actual Budget, Firefly III, Ghostfolio}
}

\pagestyle{fancy}
\fancyhf{}
\lhead{Open-Source Financial Tools}
\rhead{Desktop vs. Self-Hosted}
\cfoot{\thepage\ of \pageref{LastPage}}
\renewcommand{\headrulewidth}{0.4pt}
\renewcommand{\footrulewidth}{0pt}

\titleformat{\section}{\Large\bfseries\color{accent}}{}{0em}{}[\titlerule]
\titleformat{\subsection}{\large\bfseries\color{ink}}{}{0em}{}
\titleformat{\subsubsection}{\bfseries\color{ink}}{}{0em}{}
\setlist[itemize]{leftmargin=1.3em, itemsep=0.25em, topsep=0.25em}
\setlist[enumerate]{leftmargin=1.5em, itemsep=0.25em, topsep=0.25em}
\renewcommand{\arraystretch}{1.25}
\setlength{\parindent}{0pt}
\setlength{\parskip}{0.6em}
\setlength{\headheight}{14pt}
\setlength{\tabcolsep}{4pt}
\sloppy
\emergencystretch=3em

\newcommand{\tool}[1]{\textbf{#1}}
\newcommand{\decision}[1]{\textbf{\color{accent}#1}}
\newcommand{\smallcaps}[1]{\textsc{#1}}
\newcommand{\callout}[1]{%
  \vspace{0.2em}
  \noindent\fcolorbox{accent}{accentlight}{\parbox{0.965\linewidth}{#1}}
  \vspace{0.2em}
}

\begin{document}

\begin{center}
  {\Huge\bfseries\color{ink} Open-Source Financial Tools}\\[0.3em]
  {\Large\color{accent} Desktop and Self-Hosted Options for Budgeting, Accounting, and Portfolio Tracking}\\[0.8em]
  {\color{muted}Citation-clean handout. Source review completed May 14, 2026.}
\end{center}

\callout{\textbf{Purpose.} This handout compares credible open-source financial tools using official project documentation and established security guidance. It is intended for software selection, not financial, tax, legal, or investment advice.}

\section{Executive Summary}

Open-source financial tools can reduce vendor lock-in, improve transparency, and keep financial records under your control. The right choice depends on the financial workflow you need to operate:

\begin{itemize}
  \item Use \tool{GnuCash} when accounting correctness, double-entry ledgers, investment accounts, and small-business accounting features matter most. The project describes GnuCash as personal and small-business financial-accounting software with double-entry accounting, investment account support, reports, scheduled transactions, and import capabilities.\cite{gnucash-home,gnucash-features}
  \item Use \tool{HomeBank} or \tool{Money Manager Ex} when the goal is practical household finance management, spending analysis, budgeting, and simple net-worth visibility without running a server. HomeBank emphasizes personal accounting, import/export, templates, split categories, transfers, multiple currencies, filters, and dynamic reports.\cite{homebank-home} Money Manager Ex describes itself as free/libre, open-source, cross-platform personal finance software for tracking cash flow and everyday finances.\cite{mmex-home,mmex-docs}
  \item Use \tool{Actual Budget} when the primary job is envelope budgeting with local-first behavior, optional sync, and optional end-to-end encryption. Actual's own documentation says end-to-end encryption is optional, and its bank-sync documentation clarifies that bank-sync secrets and keys are stored separately on the server and are not covered by that encryption.\cite{actual-home,actual-sync,actual-bank-sync}
  \item Use \tool{Firefly III} when you want a self-hosted personal finance web application with budgets, categories, tags, import workflows, reports, and API-based automation. Firefly III positions itself as a self-hosted personal finance manager rather than a general small-business accounting platform.\cite{firefly-home,firefly-docs,firefly-github}
  \item Use \tool{Ghostfolio} when the core need is portfolio and net-worth tracking across assets such as cash, stocks, ETFs, and cryptocurrency. Ghostfolio describes itself as an open-source wealth management application under AGPL-3.0.\cite{ghostfolio-home,ghostfolio-github,ghostfolio-open}
\end{itemize}

\section{Corrected Source Quality}

The original source list mixed official project sites with lower-value or off-topic sources. This handout retains sources that are suitable for a technical selection document and removes sources that are weak for product claims.

\begin{tabularx}{\linewidth}{>{\RaggedRight\arraybackslash}p{0.22\linewidth} >{\RaggedRight\arraybackslash}X >{\RaggedRight\arraybackslash}p{0.26\linewidth}}
\rowcolor{headergray}\textbf{Source Type} & \textbf{Use in this handout} & \textbf{Quality Decision} \\
\toprule
Official project websites and docs & Used for feature claims, deployment model, licensing notes, and privacy statements. & Primary source. Preferred. \\
Official GitHub repositories & Used when the project website is brief or when the repository provides licensing and project-positioning evidence. & Primary source. Preferred. \\
OWASP and NIST guidance & Used for security operating practices such as MFA, secrets management, container hardening, and risk-management framing. & Authoritative security guidance. \\
SEO listicles, sponsored deployment pages, YouTube videos, broad directories, unrelated crypto-wallet comparisons & Not used for product claims. These may be useful for discovery, but they are not reliable enough for a citation-clean handout. & Excluded. \\
\bottomrule
\end{tabularx}

\section{Comparison Table}

\begingroup
\small
\rowcolors{2}{white}{headergray}
\begin{longtable}{>{\RaggedRight\arraybackslash}p{0.13\linewidth} >{\RaggedRight\arraybackslash}p{0.16\linewidth} >{\RaggedRight\arraybackslash}p{0.20\linewidth} >{\RaggedRight\arraybackslash}p{0.27\linewidth} >{\RaggedRight\arraybackslash}p{0.16\linewidth}}
\rowcolor{accentlight}\textbf{Tool} & \textbf{Deployment Model} & \textbf{Best Fit} & \textbf{Strengths} & \textbf{Watch Items} \\
\toprule
\endfirsthead
\rowcolor{accentlight}\textbf{Tool} & \textbf{Deployment Model} & \textbf{Best Fit} & \textbf{Strengths} & \textbf{Watch Items} \\
\toprule
\endhead
\tool{GnuCash} & Desktop application; local data file; GNU/Linux, BSD, Solaris, macOS, and Windows.\cite{gnucash-home} & Formal personal accounting, small-business bookkeeping, investment tracking, and double-entry workflows. & Double-entry accounting; stock, bond, and mutual-fund accounts; small-business accounting; reports and graphs; scheduled transactions; QIF/OFX/HBCI import and transaction matching.\cite{gnucash-home,gnucash-features} & More accounting-oriented learning curve. Best when the user is willing to learn accounts, ledgers, splits, and reconciliations. \\
\tool{HomeBank} & Desktop application; official downloads for Windows plus source; personal accounting software.\cite{homebank-home,homebank-downloads} & Household budgeting, spending analysis, and personal account tracking. & Import/export; templates; category split; internal transfers; multiple currencies; filters; dynamic report charts; budgeting and scheduling features.\cite{homebank-home} & Better for personal finance than formal business accounting. Multi-user or remote access is not its primary model. \\
\tool{Money Manager Ex} & Cross-platform desktop/mobile ecosystem; open-source personal finance software.\cite{mmex-home,mmex-docs} & Simple personal finance, cash-flow tracking, net-worth snapshots, and day-to-day finance organization. & Open-source and cross-platform; tracks cash flow; emphasizes simplicity and everyday use; supports common personal-finance domains such as accounts, transactions, budgets, stocks, assets, reports, import/export, and attachments.\cite{mmex-home,mmex-docs} & Lighter than GnuCash for rigorous accounting. Sync, mobile, and storage choices should be reviewed before committing sensitive records. \\
\tool{Actual Budget} & Local-first app with synchronization server; can be self-hosted.\cite{actual-home,actual-github} & Envelope budgeting and zero-based monthly planning. & Fast budgeting workflow; envelope budgeting; user-owned data; multi-device sync; optional end-to-end encryption.\cite{actual-home,actual-sync,actual-envelope} & End-to-end encryption is optional. Bank-sync API secrets and keys are not covered by that encryption, according to Actual's own bank-sync documentation.\cite{actual-bank-sync} \\
\tool{Firefly III} & Self-hosted web application.\cite{firefly-home,firefly-docs} & Power-user personal finance, categorization, reporting, imports, and automation. & Budgets, categories, tags, external import tooling, financial reports, and API-friendly operation.\cite{firefly-github,firefly-docs} & Requires server operations: patching, backups, database care, access control, and network hardening. Do not treat it as a small-business accounting substitute without validating required accounting controls. \\
\tool{Ghostfolio} & Self-hosted or hosted web application; AGPL-3.0 open-source project.\cite{ghostfolio-home,ghostfolio-open} & Portfolio, allocation, and net-worth tracking across asset classes. & Tracks net worth including cash, stocks, ETFs, and cryptocurrencies; designed as open-source wealth management software.\cite{ghostfolio-home,ghostfolio-github} & Portfolio trackers are not budgeting or accounting systems. Market-data integrations, brokerage imports, and API keys can create additional privacy exposure. \\
\bottomrule
\end{longtable}
\rowcolors{2}{}{}
\endgroup

\section{Desktop vs. Self-Hosted Decision Guidance}

\subsection{Choose a Desktop Tool When}

\begin{itemize}
  \item \decision{You want the simplest privacy boundary.} A local desktop file is easier to reason about than a public web service, reverse proxy, database, and sync server.
  \item \decision{You do not need browser-based multi-device access.} Desktop applications reduce the operational scope when a single machine or a manually synchronized file is sufficient.
  \item \decision{You need accounting discipline.} GnuCash is the strongest fit in this set for double-entry accounting and small-business accounting workflows.\cite{gnucash-home,gnucash-features}
  \item \decision{You prefer minimal infrastructure.} No container runtime, database server, TLS certificate, identity provider, or uptime monitoring is required for the basic desktop model.
\end{itemize}

\subsection{Choose a Self-Hosted Web Tool When}

\begin{itemize}
  \item \decision{You need access from multiple devices.} Web applications such as Actual Budget, Firefly III, and Ghostfolio fit browser-based access and synchronized workflows.\cite{actual-home,firefly-docs,ghostfolio-home}
  \item \decision{You want automation and integration.} Firefly III's API and import-oriented ecosystem are better aligned to data-ingestion and reporting workflows than a purely local desktop file.\cite{firefly-github,firefly-docs}
  \item \decision{You can operate the server responsibly.} Self-hosting shifts responsibility for patching, backups, TLS, identity, secrets, and monitoring to you.
  \item \decision{You accept a larger attack surface.} A self-hosted finance app can be private, but only if the deployment is hardened and maintained.
\end{itemize}

\callout{\textbf{Selection rule.} Start with the workflow, not the tool. If the workflow is formal accounting, start with GnuCash. If the workflow is monthly envelope budgeting, start with Actual Budget. If the workflow is portfolio/net-worth visibility, start with Ghostfolio. If the workflow is self-hosted personal finance reporting and automation, evaluate Firefly III.}

\section{Privacy and Security Considerations}

Financial records are sensitive operational data. Open source does not automatically make a finance system secure; it makes the implementation inspectable and gives the operator more control. The operating model still determines the risk.

\subsection{Data Location and Control}

\begin{itemize}
  \item Document where data lives: local files, application database, backup storage, mobile device, sync server, bank-sync provider, and exported CSV/OFX/QIF files.
  \item Encrypt backups before moving them to cloud storage or removable media.
  \item Treat exported statements and transaction files as sensitive records. They often contain account names, transaction descriptions, merchants, balances, and timestamps.
  \item For Actual Budget, distinguish budget-data encryption from bank-sync token handling. Actual's docs state that bank-sync secrets and keys are stored separately on the server and are not covered by end-to-end encryption.\cite{actual-bank-sync}
\end{itemize}

\subsection{Access Control}

\begin{itemize}
  \item Use strong authentication and enable MFA wherever the application or fronting identity provider supports it. OWASP identifies MFA as a strong defense against password-related attacks.\cite{owasp-mfa,owasp-auth}
  \item Separate administrator accounts from normal user accounts.
  \item Do not reuse banking, brokerage, or finance-tool passwords across services.
  \item If exposing a web application remotely, prefer VPN or trusted network access over direct public exposure.
\end{itemize}

\subsection{Secrets and Bank Integrations}

\begin{itemize}
  \item Keep API keys, bank-sync tokens, database passwords, SMTP passwords, and OAuth secrets out of source control.
  \item Apply least privilege to secrets: only the component that needs a credential should be able to read it. OWASP's secrets-management guidance emphasizes fine-grained access control and least privilege for secrets.\cite{owasp-secrets}
  \item Rotate credentials after suspected compromise, staff/device changes, or accidental exposure.
  \item Prefer read-only bank or brokerage connections where available.
\end{itemize}

\subsection{Container and Server Hardening for Self-Hosted Apps}

\begin{itemize}
  \item Patch the application, base images, database, reverse proxy, and host operating system on a defined cadence.
  \item Avoid privileged containers, avoid mounting the Docker socket into application containers, and avoid running services as root when the image supports non-root execution. OWASP's Docker Security Cheat Sheet highlights common container misconfigurations and hardening practices.\cite{owasp-docker}
  \item Bind services to private networks where possible; expose only the reverse proxy or VPN endpoint.
  \item Use TLS for browser access, even on private networks, when practical.
  \item Monitor logs for failed logins, unexpected admin actions, backup failures, and import failures.
\end{itemize}

\subsection{Backup and Recovery}

\begin{itemize}
  \item Back up both application data and configuration. For self-hosted apps, this usually includes database dumps, uploaded files, environment configuration, and compose/Kubernetes manifests.
  \item Test restore procedures. A backup that has never been restored is an assumption, not a control.
  \item Keep at least one offline or immutable backup to reduce ransomware and accidental-deletion exposure.
  \item Align operating practices to a basic risk-management lifecycle: govern, identify, protect, detect, respond, and recover. NIST CSF 2.0 uses these six functions to organize cybersecurity outcomes.\cite{nist-csf}
\end{itemize}

\section{Recommended Shortlist by Use Case}

\begin{tabularx}{\linewidth}{>{\RaggedRight\arraybackslash}p{0.28\linewidth} >{\RaggedRight\arraybackslash}p{0.27\linewidth} >{\RaggedRight\arraybackslash}X}
\rowcolor{headergray}\textbf{Use Case} & \textbf{Primary Recommendation} & \textbf{Rationale} \\
\toprule
Personal double-entry accounting & GnuCash & Most accounting-complete option in this set; strong ledger model and reporting. \\
Small-business bookkeeping on a desktop & GnuCash & Explicitly positioned for personal and small-business financial accounting. \\
Simple household finance tracking & HomeBank or Money Manager Ex & Lower operational complexity than self-hosting; easier personal-finance workflows. \\
Envelope budgeting & Actual Budget & Purpose-built around envelope budgeting with local-first behavior and optional sync. \\
Self-hosted personal finance dashboard & Firefly III & Strongest fit for users comfortable maintaining a web app, database, imports, and reports. \\
Investment allocation and net-worth tracking & Ghostfolio & Designed for wealth management and portfolio visibility rather than ledger accounting. \\
Privacy-minimal setup & Desktop tool with encrypted local backups & Avoids the server, reverse proxy, database, and internet-exposure problem entirely. \\
Multi-device access with control & Self-hosted Actual, Firefly III, or Ghostfolio behind VPN/SSO & Best when remote access is required and the operator can maintain the infrastructure. \\
\bottomrule
\end{tabularx}

\section{Implementation Checklist}

\subsection{Before Choosing a Tool}

\begin{enumerate}
  \item Define the primary workflow: accounting, budgeting, expense tracking, net-worth tracking, or portfolio analysis.
  \item Decide whether multi-device browser access is required or merely convenient.
  \item List import sources: CSV, QIF, OFX, bank sync, brokerage exports, crypto exchange exports, or manual entry.
  \item Decide who needs access and whether shared access is required.
  \item Define the backup target, encryption method, retention period, and restore test schedule.
\end{enumerate}

\subsection{Before Self-Hosting}

\begin{enumerate}
  \item Place the app behind VPN, private networking, or a hardened reverse proxy.
  \item Enable MFA through the app or identity provider when available.
  \item Store secrets outside source control and restrict access.
  \item Run containers with least privilege where possible.
  \item Automate backups and test restoration.
  \item Subscribe to project release notes or repository security advisories.
  \item Document the recovery procedure in the same repository as the deployment manifests.
\end{enumerate}

\section{Bottom Line}

For most users, \tool{desktop tools} offer the best privacy-to-effort ratio. For power users, \tool{self-hosted tools} provide better automation and cross-device access but require real operational discipline. A practical stack might use \tool{GnuCash} for formal accounting, \tool{Actual Budget} for household envelope budgeting, and \tool{Ghostfolio} for portfolio visibility, while avoiding unnecessary overlap and duplicate data entry.

\newpage
\begin{thebibliography}{99}

\bibitem{gnucash-home} GnuCash Project. \emph{GnuCash: Free Accounting Software}. \url{https://gnucash.org/}

\bibitem{gnucash-features} GnuCash Project. \emph{GnuCash Features}. \url{https://gnucash.org/features.phtml}

\bibitem{homebank-home} HomeBank. \emph{HomeBank: Free personal finance software}. \url{https://www.gethomebank.org/}

\bibitem{homebank-downloads} HomeBank. \emph{HomeBank Downloads}. \url{https://www.gethomebank.org/en/downloads.php}

\bibitem{mmex-home} Money Manager Ex. \emph{Money Manager Ex: Free, easy-to-use personal finance software}. \url{https://moneymanagerex.org/}

\bibitem{mmex-docs} Money Manager Ex. \emph{Money Manager Ex User Manual and Features}. \url{https://moneymanagerex.org/moneymanagerex/}

\bibitem{actual-home} Actual Budget. \emph{Actual Budget: Your Finances - made simple}. \url{https://actualbudget.org/}

\bibitem{actual-github} Actual Budget. \emph{actualbudget/actual GitHub Repository}. \url{https://github.com/actualbudget/actual}

\bibitem{actual-sync} Actual Budget Documentation. \emph{Syncing Across Devices}. \url{https://actualbudget.org/docs/getting-started/sync}

\bibitem{actual-envelope} Actual Budget Documentation. \emph{Envelope Budgeting}. \url{https://actualbudget.org/docs/getting-started/envelope-budgeting/}

\bibitem{actual-bank-sync} Actual Budget Documentation. \emph{Connecting Your Bank}. \url{https://actualbudget.org/docs/advanced/bank-sync}

\bibitem{firefly-home} Firefly III. \emph{A free and open source personal finance manager}. \url{https://www.firefly-iii.org/}

\bibitem{firefly-docs} Firefly III Documentation. \emph{Introduction and Features}. \url{https://docs.firefly-iii.org/explanation/firefly-iii/about/introduction/}

\bibitem{firefly-github} Firefly III. \emph{firefly-iii/firefly-iii GitHub Repository}. \url{https://github.com/firefly-iii/firefly-iii}

\bibitem{ghostfolio-home} Ghostfolio. \emph{Open Source Wealth Management Software}. \url{https://ghostfol.io/}

\bibitem{ghostfolio-github} Ghostfolio. \emph{ghostfolio/ghostfolio GitHub Repository}. \url{https://github.com/ghostfolio/ghostfolio}

\bibitem{ghostfolio-open} Ghostfolio. \emph{Open Startup}. \url{https://ghostfol.io/en/open}

\bibitem{owasp-docker} OWASP Cheat Sheet Series. \emph{Docker Security Cheat Sheet}. \url{https://cheatsheetseries.owasp.org/cheatsheets/Docker_Security_Cheat_Sheet.html}

\bibitem{owasp-mfa} OWASP Cheat Sheet Series. \emph{Multifactor Authentication Cheat Sheet}. \url{https://cheatsheetseries.owasp.org/cheatsheets/Multifactor_Authentication_Cheat_Sheet.html}

\bibitem{owasp-auth} OWASP Cheat Sheet Series. \emph{Authentication Cheat Sheet}. \url{https://cheatsheetseries.owasp.org/cheatsheets/Authentication_Cheat_Sheet.html}

\bibitem{owasp-secrets} OWASP Cheat Sheet Series. \emph{Secrets Management Cheat Sheet}. \url{https://cheatsheetseries.owasp.org/cheatsheets/Secrets_Management_Cheat_Sheet.html}

\bibitem{nist-csf} National Institute of Standards and Technology. \emph{The NIST Cybersecurity Framework (CSF) 2.0}. NIST CSWP 29, February 26, 2024. \url{https://nvlpubs.nist.gov/nistpubs/CSWP/NIST.CSWP.29.pdf}

\end{thebibliography}

\end{document}
